\section*{Chapitre \ref{ch:g6:inside-the-earth} --- L'intérieur de la Terre}

\begin{solution}{exo:g6:inside-the-earth:1}
La \omterm{def:g6:inside-the-earth:layers}{croûte}, le \omterm{def:g6:inside-the-earth:layers}{manteau}, le \omterm{def:g6:inside-the-earth:layers}{noyau externe}, le \omterm{def:g6:inside-the-earth:layers}{noyau interne}. Le \omterm{def:g6:inside-the-earth:layers}{noyau
externe} est liquide. Le \omterm{def:g6:inside-the-earth:layers}{noyau externe} et le \omterm{def:g6:inside-the-earth:layers}{noyau interne} sont faits de
métal (surtout du fer, avec un peu de nickel)~; la \omterm{def:g6:inside-the-earth:layers}{croûte} et le \omterm{def:g6:inside-the-earth:layers}{manteau}
sont faits de \omterm{def:g3:stones-pebbles-sand:rock}{roche}.
\end{solution}

\begin{solution}{exo:g6:inside-the-earth:2}
(a) Faux~: la \omterm{def:g6:inside-the-earth:layers}{croûte} est plus épaisse sous les \omterm{def:g3:continents-and-oceans:continent}{continents} (environ 30~km)
que sous les \omterm{def:g2:water-on-earth:ocean}{océans} (environ 5~km). (b) Faux~: le \omterm{def:g6:inside-the-earth:layers}{noyau} est fait de métal,
un alliage de fer et de nickel. (c) Vrai. (d) Vrai~: 12\,262~m, environ
12.3~km.
\end{solution}

\begin{solution}{exo:g6:inside-the-earth:3}
(a) À 2891~km, environ 2900~km. (b) À 5150~km. (c) $5150 - 2891 = 2259$~:
environ 2260~km.
\end{solution}

\begin{solution}{exo:g6:inside-the-earth:4}
«~La coquille est comme la \underline{\omterm{def:g6:inside-the-earth:layers}{croûte}}, le blanc est comme le
\underline{\omterm{def:g6:inside-the-earth:layers}{manteau}} et le jaune est comme le \underline{\omterm{def:g6:inside-the-earth:layers}{noyau}}. Mais
contrairement au jaune, le \underline{\omterm{def:g6:inside-the-earth:layers}{noyau}} a deux parties~: un
\underline{\omterm{def:g6:inside-the-earth:layers}{noyau externe}} liquide et un \underline{\omterm{def:g6:inside-the-earth:layers}{noyau interne}}
solide.~»
\end{solution}

\begin{solution}{exo:g6:inside-the-earth:5}
$8.849 \div 6371 \approx 0.00139$, et $0.00139 \times 100 \approx 0.14$~:
l'Everest ne représente qu'environ \textbf{0.14~\%} du rayon. Vue de
loin, la Terre est une boule très lisse~: sa plus haute \omterm{def:g1:landscapes-around-us:landform}{montagne} n'est
qu'une minuscule bosse.
\end{solution}

\begin{solution}{exo:g6:inside-the-earth:6}
(a) $30 \div 5 = 6$~: six fois plus épaisse. (b) $100 \div 5 = 20$~: vingt
fois plus épaisse.
\end{solution}

\begin{solution}{exo:g6:inside-the-earth:7}
(a) $12 + 1 \times 25 = 37~^\circ$C. (b) $12 + 1.5 \times 25 = 12 + 37.5 =
49.5~^\circ$C. (c) $12 + 2 \times 25 = 62~^\circ$C.
\end{solution}

\begin{solution}{exo:g6:inside-the-earth:8}
La hausse est de $74 - 14 = 60$ degrés sur 2~km, donc $60 \div 2 =
\textbf{30}$~degrés par kilomètre. Oui~: 30 est compris entre la moyenne
de l'est (25) et celle de l'ouest (34).
\end{solution}

\begin{solution}{exo:g6:inside-the-earth:9}
(a) Juste au-dessus du \omterm{def:g6:inside-the-earth:layers}{noyau}, environ 13.7~km par seconde~; juste sous son
sommet, environ 8.1~km par seconde~: les ondes P ralentissent brusquement.
(b) Entre 2891~km et 5150~km, dans le \omterm{def:g6:inside-the-earth:layers}{noyau externe}. (c) Les ondes S ne
peuvent pas traverser un liquide~: le \omterm{def:g6:inside-the-earth:layers}{noyau externe} est liquide.
\end{solution}

\begin{solution}{exo:g6:inside-the-earth:10}
(a) $12.3 \div 30 = 0.41$~: environ \textbf{41~\%} d'une \omterm{def:g6:inside-the-earth:layers}{croûte}
continentale moyenne. (b) Non~: sous un \omterm{def:g3:continents-and-oceans:continent}{continent}, le \omterm{def:g6:inside-the-earth:layers}{manteau} commence
sous la \omterm{def:g6:inside-the-earth:layers}{croûte}, à environ 30~km de profondeur en moyenne~; le forage s'est
arrêté bien au-dessus, dans la \omterm{def:g6:inside-the-earth:layers}{croûte}.
\end{solution}

\begin{solution}{exo:g6:inside-the-earth:11}
(a) $29 \times 6371 = 184\,759$, et $15 + 184\,759 = 184\,774$~: environ
185\,000~degrés. (b) Le laboratoire donne environ 6000~degrés à la limite
du \omterm{def:g6:inside-the-earth:layers}{noyau interne}~: le nombre de Léo est environ 30~fois trop grand
($184\,774 \div 6000 \approx 31$). (c) Le rythme de 29~degrés par
kilomètre n'a été mesuré que dans les 2~premiers km environ. Léo l'utilise
des milliers de kilomètres au-delà des mesures~; en profondeur, la
température monte beaucoup plus lentement. Un rythme ne peut être utilisé
que là où il a été mesuré.
\end{solution}

\begin{solution}{exo:g6:inside-the-earth:12}
(a) Si la \omterm{def:g6:inside-the-earth:layers}{croûte} est plus légère que la moyenne, les matériaux plus
profonds doivent avoir une masse volumique supérieure à 5.5, pour faire
monter la moyenne. (b) Dans le \omterm{def:g6:inside-the-earth:layers}{noyau externe} et le \omterm{def:g6:inside-the-earth:layers}{noyau interne} (d'environ
9.9 à environ 13~grammes par centimètre cube)~; le \omterm{def:g6:inside-the-earth:layers}{manteau} n'atteint tout
juste 5.6 qu'à sa base. (c) Les matériaux plus denses coulent sous les
plus légers, comme une pierre coule dans l'eau~: le métal lourd a fini au
milieu et les \omterm{def:g3:stones-pebbles-sand:rock}{roches} plus légères en haut.
\end{solution}

\begin{solution}{pb:g6:inside-the-earth:1}
\textbf{1.} $2 \times 6371 = 12\,742$~km. La maquette a 2~m de diamètre,
soit 2000~mm.

\textbf{2.} $12\,742 \div 2000 \approx 6.37$~: 1~mm de la maquette
représente environ 6.37~km de la vraie Terre.

\textbf{3.} $8.849 \div 6.371 \approx 1.39$~: l'Everest mesure environ
\textbf{1.4~mm} de haut sur la maquette, un grain de \omterm{def:g3:stones-pebbles-sand:sand}{sable}.

\textbf{4.} $12.3 \div 6.371 \approx 1.9$~mm. Le centre de la maquette est
à 1000~mm (1~m) sous la surface~: le forage est environ 520~fois plus
court ($1000 \div 1.9 \approx 520$).

\textbf{5.} $2891 \div 6.371 \approx 453.8$~: le haut du \omterm{def:g6:inside-the-earth:layers}{noyau externe} se
peint à environ \textbf{454~mm} sous la surface.

\textbf{6.} $5150 \div 6.371 \approx 808.3$~: environ 808~mm sous la
surface. Le rayon du \omterm{def:g6:inside-the-earth:layers}{noyau interne} sur la maquette est de
$1000 - 808 = 192$~mm.

\textbf{7.} Le \omterm{def:g6:inside-the-earth:layers}{manteau}~: d'environ 5~mm à 454~mm, soit environ 449~mm. Le
\omterm{def:g6:inside-the-earth:layers}{noyau externe} occupe $808 - 454 = 354$~mm et le \omterm{def:g6:inside-the-earth:layers}{noyau interne} 192~mm.

\textbf{8.} $30 \div 6.371 \approx 4.7$~mm~: une bande fine, que l'on peut
peindre avec un pinceau fin.

\textbf{9.} $15 + 2 \times 29 = 73~^\circ$C. Sur la maquette, 2~km se
trouvent à $2 \div 6.371 \approx 0.31$~mm sous la surface~: à peine
l'épaisseur de trois feuilles de papier.

\textbf{10.} $4 \div 6.371 \approx 0.63$~mm, contre environ 1.9~mm pour le
forage~: la mine est moins profonde, d'environ $1.9 - 0.6 = 1.3$~mm.

\textbf{11.} Qu'une substance fonde ou non dépend de la température et de
la pression. Le \omterm{def:g6:inside-the-earth:layers}{noyau interne} est le plus comprimé de tous, par environ
3.3~millions de fois la pression de l'air~: le fer comprimé ne fond qu'à
une température plus élevée, si bien que le fer du \omterm{def:g6:inside-the-earth:layers}{noyau interne} reste
solide alors qu'il est plus chaud que le \omterm{def:g6:inside-the-earth:layers}{noyau externe}.

\textbf{12.} $5 \div 6.371 \approx 0.785$~: la \omterm{def:g6:inside-the-earth:layers}{croûte} océanique ne doit
être peinte que sur \textbf{0.78~mm} d'épaisseur sur la maquette de
2~mètres, l'épaisseur d'environ huit feuilles de papier
($0.78 \div 0.1 \approx 8$).
\end{solution}
